\chapter{Contexto y estado de la cuestión}

El oscilador armónico cuántico (OAC) es uno de los sistemas modelo más importantes en la mecánica cuántica \citep{schrodinger1926a, bornheisenberg1926, dirac1930}, siendo el análogo cuántico del oscilador armónico clásico descrito a partir de la dinámica newtoniana \citep{newton1687}. Además, es uno de los pocos sistemas mecánico-cuánticos que admiten una solución exacta y analíticamente tratable \citep{schrodinger1926a}. Este sistema simple permite explicar una amplia variedad de fenómenos físicos, desde las vibraciones moleculares y modos normales en sólidos \citep{einstein1907,debye1912} hasta la cuantización del campo electromagnético y el comportamiento cuántico de la luz \citep{dirac1927}.

\section{El problema de la radiación del cuerpo negro}

En el siglo XIX se establece definitivamente la teoría electromagnética de la luz \citep{maxwell1865} y se investiga su aplicación al estudio de la emisión de radiación de los átomos \citep{bohr1913}. Esta teoría se podía aplicar a casi todos los experimentos \citep{rayleigh1900}, existían algunas excepciones donde aparecían ciertos desajustes \citep{jeans1905}. Al estudiar la emisión de radiación producida por un cuerpo en equilibrio térmico \citep{wien1896}, la teoría clásica no era capaz de explicar su comportamiento a altas frecuencias de la emisión \citep{rayleigh1900, jeans1905}.

Las expresiones clásicas propuestas por \cite{rayleigh1900} y desarrolladas por \cite{jeans1905} describían correctamente la radiación a bajas frecuencias. Sin embargo, a altas frecuencias la densidad de energía predicha divergía en el rango ultravioleta, en desacuerdo con los resultados experimentales \citep{wien1896}. Este problema, conocido como la ``catástrofe ultravioleta'', evidenció una limitación de la teoría clásica y motivó la introducción de la hipótesis de cuantización propuesta por \cite{planck1901}

Una solución matemática al problema de la radiación del cuerpo negro fue propuesta por \cite{planck1900, planck1901}. Inicialmente presentada como un recurso formal para ajustar los datos experimentales \citep{wien1896}, esta expresión reproducía con gran precisión el espectro observado, lo que motivó posteriormente un análisis más profundo de sus implicancias físicas.

\cite{planck1901} modeló la emisión de radiación mediante un oscilador armónico cargado y postuló que la energía se intercambiaba en paquetes discretos o ``cuantos''. Según esta hipótesis, la energía del oscilador solo podía tomar valores  $E_n = nh\nu$, donde $n$ es un número entero positivo \citep{planck1900, planck1901}. Esta idea rompió con la física clásica continua y abrió el camino al desarrollo de la teoría cuántica \citep{einstein1907, bornheisenberg1926}.

\cite{einstein1907} profundizó la propuesta de \cite{planck1901}. En ese mismo trabajo de 1907, se encontraba analizando otro problema aparentemente no relacionado: ¿por qué el calor específico de los sólidos disminuye drásticamente cuando se enfrían? La ley clásica de Dulong y Petit \citep{dulong1819} establecía que el calor específico debía ser constante, pero los experimentos mostraban claramente que no era así \citep{dulong1819}. \cite{einstein1907} planteó que, si los átomos en un sólido vibran como osciladores armónicos y su energía está cuantizada según la hipótesis de \cite{planck1901}, entonces a temperatura $T$ la energía promedio de cada oscilador viene dada por
\vspace{-10pt}
\begin{equation}
	\langle E \rangle = \frac{h\nu}{e^{h\nu/kT} - 1}.
\end{equation}

Esta fórmula explicaba exactamente lo que se observaba en los experimentos \citep{einstein1907}. Fue la primera vez que alguien usó el oscilador armónico cuántico para explicar un fenómeno físico real \citep{einstein1907}, no solo la radiación \citep{planck1901}. \citep{debye1912} mejoró el modelo en 1912 considerando que los átomos pueden vibrar a diferentes frecuencias, pero la idea central de Einstein quedó intacta \citep{einstein1907, debye1912}.


\subsection*{Dos caminos hacia la misma teoría}

En la década de 1920 se desarrollaron nuevas formulaciones matemáticas de la teoría cuántica iniciada por Planck, Einstein y Debye \citep{planck1901, einstein1907, debye1912}. En este contexto, Schrödinger propuso en 1926 describir las partículas como ondas y formular su dinámica como un problema de autovalores \citep{schrodinger1926a, schrodinger1926b}.

\vspace{-20pt}
\begin{equation}
	\hat{H}\psi = E\psi
\end{equation}donde $\hat{H}$ es el operador hamiltoniano y $E$ el valor propio de energía asociado al estado $\psi$ \citep{schrodinger1926a, schrodinger1926b}. En el caso del oscilador armónico, la exigencia de soluciones normalizables conduce a un espectro discreto de energías \citep{schrodinger1926a}. Así, la cuantización surge naturalmente de la ecuación de Schrödinger, y el comportamiento clásico puede recuperarse como un caso límite \citep{ehrenfest1927}.

\cite{heisenberg1925} había tomado un camino totalmente distinto. Planteó trabajar solo con las cantidades que sí son observables. No usa trayectorias ni posiciones exactas \citep{heisenberg1927}. Estas frecuencias están relacionadas con las energías mediante la denominada relación de frecuencia de Bohr \citep{bohr1913}:
\vspace{-10pt}
\begin{equation}
	\nu_{mn} = \frac{E_m - E_n}{h}
\end{equation}

En este enfoque estas magnitudes obedecen reglas de combinación no conmutativas, lo que implica que el orden de las operaciones es relevante en su tratamiento matemático \citep{heisenberg1925}. Es decir, $AB$ no es lo mismo que $BA$. \cite{born1925} formalizaron esto matemáticamente, y junto con Heisenberg \cite{bornheisenberg1926} demostraron que la relación fundamental es:
\vspace{-10pt}
\begin{equation}
	[\hat{q}, \hat{p}] = i\hbar
\end{equation}

A partir de esta relación \citep{bornheisenberg1926} se puede derivar el espectro de energías del oscilador sin tener que resolver ninguna ecuación diferencial, solo usando métodos algebraicos \citep{dirac1930}. De este proceso surgen los operadores de creación y aniquilación que simplifican enormemente todos los cálculos asociados al oscilador armónico \citep{dirac1930}.

\section{De partículas a campos}

\cite{dirac1927} extendió el principio de cuantización al campo electromagnético, describiéndolo como un conjunto de modos normales equivalentes a osciladores armónicos cuánticos. Este enfoque permitió explicar procesos de emisión y absorción de radiación y fue posteriormente sistematizado en su tratado \citep{dirac1930}.

Esto no solo era una analogía matemática. \cite{dirac1927} pudo explicar con este modelo cosas que antes requerían suposiciones adicionales, como por qué los átomos emiten luz espontáneamente incluso sin ningún estímulo externo \citep{dirac1927}. A partir de este momento, el oscilador armónico \citep{dirac1930} dejó de ser un modelo pedagógico y se convirtió en el lenguaje natural para describir diversos fenómenos cuánticos \citep{dirac1930, bornheisenberg1926}. 

Todo lo anterior explica por qué el oscilador armónico \citep{schrodinger1926a, dirac1930} se convirtió en un sistema fundamental sobre el que se construye gran parte de la física cuántica moderna \citep{dirac1930, bornheisenberg1926}. Realizar el cálculo de sus niveles de energía y sus funciones de onda ---objetivo de este trabajo--- no será un ejercicio matemático \citep{dirac1930}. Nos permitirá entender el mecanismo básico \citep{dirac1930} que aparece en problemas tan diversos de la mecánica cuántica como: vibraciones moleculares \citep{einstein1907}, modos normales en sólidos \citep{debye1912}, y la cuantización del campo electromagnético \citep{dirac1927}, entre otros.